\section{Conclusion}\label{sec:conclusion}

A GPU-accelerated D3Q19 MRT lattice Boltzmann solver coupled with
Bayesian optimization was applied to the design of gyroid TPMS
catalyst supports.  The main findings are:

\begin{enumerate}
  \item \textbf{Efficient design-space exploration:}
    100 BO evaluations (85 feasible) produced 50 Pareto-optimal
    designs in the $S_v$--$\Delta P$ space, spanning
    $a = \SIrange{3}{12}{\milli\meter}$ and $t \in (-0.5, 0.5)$.

  \item \textbf{Design selection:}
    Five representative designs were identified.
    Design~E ($a = \SI{5.17}{\milli\meter}$, $\varepsilon = 0.44$,
    $S_v = \SI{2252}{\per\meter}$, $\Delta P = \SI{20}{\pascal}$)
    offers the best balance of catalytic surface area, flow resistance,
    and manufacturability.

  \item \textbf{Forchheimer characterization:}
    Inertial coefficients $\beta \sim \SIrange{2000}{2900}{\per\meter}$
    were reliably extracted for four of five designs ($R^2 > 0.96$).
    The high-permeability design~B exceeded the solver's
    compressibility limit at elevated body forces.

  \item \textbf{Deterministic reproducibility:}
    The GPU-based solver demonstrated CV $< 10^{-13}\%$
    across triplicate runs, ensuring numerical consistency.
\end{enumerate}

Future work will address grid-independence studies at finer resolutions,
experimental validation with 3D-printed ceramic gyroid specimens,
and extension of the BO framework to include thermal performance metrics.

% Acknowledgements moved to main.tex
